% results.tex

We evaluated SA on six protein families with different sizes, functions, and PCA-space
geometries (Table~\ref{tab:cross-family}). Five were Pfam seed families: Kunitz
(PF00014, $K=99$), SH3 (PF00018, $K=55$), WW (PF00397, $K=420$), Homeobox
(PF00046, $K=136$), and Forkhead (PF00250, $K=246$). We used these families to test
hard curation and multiplicity-weighted conditioning. Data-selected single-position markers
defined their designated subsets. Prior literature motivated the residue classes where
available, but the markers are not direct activity labels. We also studied the
$\omega$-conotoxin O-superfamily ($K=74$, curated from SwissProt) as a peptide example
associated with the Cav2.2 calcium channel.

For each family, we split the sequences into designated and background subsets. We then built
the conditioned memory matrices and generated sequences at the entropy-crossover inverse
temperature $\beta^{*}$. The separation index $S$ measures PCA separation between the two
subsets. It ranged from \WWSeparation{} for WW to \HomeoboxSeparation{} for Homeobox. This
range allowed us to compare conditioning across different PCA geometries.

\paragraph{Hard curation and marker retention.}
We restricted the memory matrix to a K/R-at-P1-positive subset of Kunitz domains
(PF00014; $K=99$ sequences, $L=53$ aligned positions) and tested whether the selected
sequence marker transfers to generated sequences (Fig.~\ref{fig:scaling}). Kunitz domains are
serine protease inhibitors. The P1 residue at the primary contact site influences target
specificity.
Lys or Arg at P1 is associated with trypsin specificity~\cite{laskowskiProteinInhibitors1980},
but is not by itself an activity measurement for every sequence. We designated the 32 sequences
with K/R at P1 as the marker-positive subset ($K_{\mathrm{des}} = 32$) and the remaining 67 as
marker-negative background ($K_{\mathrm{bg}} = 67$). We built a separate memory matrix for
each group and selected $\beta^{*}$ from the entropy-crossover onset. Each condition produced
930 sequences from 30 chains run for 5{,}000 steps and thinned after burn-in.
All designated-subset-conditioned sequences had K/R at P1. No background-conditioned
sequence did. Full-family generation produced 38\% K/R, compared with 32\% in the input
family. Complete K/R retention is expected because the designated input was selected on this
marker. It is not independent evidence of trypsin inhibition.
Sequence-level inspection showed that the top K/R-positive-subset-conditioned sequences
preserved all 11 highly conserved positions, all six cysteines, and K/R at P1. They contained
19--26 substitutions, mainly at variable sites. Per-position Shannon entropy correlated with
the stored MSA for both full-family generation ($r = 0.988$) and K/R-positive-subset
conditioning ($r = 0.932$). The largest change occurred at P1, where conditioning restricted
the residues to K/R (Fig.~\ref{fig:sequence-analysis}).

We next varied the designated input size. We sampled
$K_{\mathrm{des}} \in \{3, 5, 8, 10, 15, 20, 25, 32\}$ Kunitz K/R-positive sequences and
used three replicates per size. P1 marker retention was 100\% at every size, including
$K_{\mathrm{des}}=3$. Pairwise sequence diversity increased from 0.35 at
$K_{\mathrm{des}}=3$ to 0.50 at $K_{\mathrm{des}}=32$. Over the same range, KL divergence
decreased from 0.070 to 0.009 (Fig.~\ref{fig:scaling}). Smaller input sets define a narrower
PCA subspace and therefore produced less diverse libraries. These results show marker
retention from small designated inputs, with diversity increasing as more inputs were added.

\paragraph{Multiplicity-weighted conditioning and the calibration gap.}
We tested multiplicity-weighted generation on the Kunitz family ($K = 99$,
$K_{\mathrm{des}} = 32$ K/R-positive sequences, $K_{\mathrm{bg}} = 67$ background),
sweeping $\rho$ from 1 (standard SA, no bias) to 500 in the canonical replicated
analysis (Fig.~\ref{fig:phase-transition}, Table~\ref{tab:per-family-rho}). For each $\rho$, we
computed the multiplicity vector $\vr$. We then found $\beta^{*}(\vr)$ from the
entropy-crossover onset of the weighted attention distribution. We ran five independent
replicates of 20 chains, with 620
retained samples per replicate. At each $\rho$, we recorded the mean attention weight on
designated patterns, $\bar{a}_{\mathrm{des}}$. We also recorded the decoded K/R fraction at
P1, $f_{\mathrm{obs}}$. Across all six canonical replicated sweeps, the largest absolute
attention deviation was \AttnMaxDevPts{} percentage points
(\AttnMaxDevFamily{}, $\rho=\AttnMaxDevRho$), showing close but not exact finite-sample
agreement with the target mixture fraction. However, the decoded marker fraction told a
different story. For Kunitz, P1~K/R reached
\KunitzFobsFiveHundred{} $\pm$ \KunitzFobsFiveHundredSD{} at $\rho=500$, far short of the
$f_{\mathrm{eff}} = 0.996$
that the attention weights achieved. This discrepancy between latent designated mass and the
decoded marker marginal is the calibration gap. As
Eq.~\ref{eq:decoder-pushforward} shows, it reflects both false negatives and
false positives under the complete latent-to-sequence map. Sequence diversity decreased
from 0.56 to 0.51 over this range. KL divergence increased from 0.007 to 0.017. The entropy crossover
$\beta^{*}$ shifted from 4.4 at $\rho=1$ ($K_{\mathrm{eff}} = 99$) to 9.3 at
$\rho=1{,}000$ ($K_{\mathrm{eff}} = 32.1$). The entropy curves shifted to the right as
$\rho$ increased (Fig.~\ref{fig:phase-transition}).

To examine this gap, we measured three layers of the signal chain
(Eq.~\ref{eq:gap-decomp}) over 17 values of $\rho$ from 1 to 1{,}000, with five replicates
per value. The measurements were designated attention mass
$\bar{a}_{\mathrm{des}}$, continuous reconstructed K/R probability
$f_{\mathrm{soft}}$, and decoded K/R fraction $f_{\mathrm{obs}}$. The attention gap remained
small. The PCA term was the largest component. The value of $f_{\mathrm{soft}}$ stayed near
0.27 and varied by less than 0.04 over the full range of $\rho$. The argmax term was negative
($\approx-0.25$), so discrete decoding recovered some signal lost during PCA reconstruction.
Even with 99.8\% of the attention on designated patterns, the continuous reconstruction did
not resolve K from G at P1. This variation was not well represented by the leading
$d=80$ principal components. The Kunitz split had separation index
$S=\KunitzSeparation$, which indicates partial PCA overlap.

\paragraph{Exact-equilibrium benchmark.}
The closed-form Gaussian-mixture target permits independent equilibrium sampling without
a Markov chain. We compared 1{,}640 exact draws with 1{,}640 retained ULA states at
$\rho\in\{1,10,500\}$, using the same $\beta^{*}$ and decoder in each comparison. At
$\rho=500$, the exact designated-component probability was
$f_{\mathrm{eff}}=0.996$, while the decoded P1~K/R fraction was 0.604 under exact
sampling and 0.597 under ULA. The amino acid composition divergence
$D_{\mathrm{KL}}(\text{exact} \parallel \text{ULA})$ between the two decoded libraries
was $1.19\times10^{-4}$. Across all three conditions, exact
and ULA MAP-designated basin occupancies differed by at most 0.026, and the decoded
P1~K/R fractions differed by at most 0.016. The full comparison is provided in the
Supporting Information. These results show that the calibration gap persists at
exact equilibrium, indicating that finite-step ULA discretization and mixing error were not
the main source of the gap.

\paragraph{Cross-family marker analysis and the separation index.}
We repeated the multiplicity analysis on four additional Pfam families. The SH3 split
(PF00018; $K=55$) used a data-selected Trp marker. The WW split
(PF00397; $K=420$) used a data-selected middle-third residue. The Homeobox split
(PF00046; $K=136$) used a data-selected Gln marker. The Forkhead split
(PF00250; $K=246$) used a data-selected His/Asn marker. Together with Kunitz, these
families had distinct PCA-space geometries
(Table~\ref{tab:cross-family}, Table~\ref{tab:per-family-rho}). Their calibration gaps varied
inversely with separation overall, but not monotonically
(Fig.~\ref{fig:separation-vs-gap}). An exploratory unweighted fit across the five Pfam
families yielded
$\Delta \approx \RelationIntercept{} \RelationSlope\,S$ with $R^{2}=\RelationRsq$.
Leave-one-family-out sensitivity was substantial: slopes ranged from
\RelationLooSlopeMin{} to \RelationLooSlopeMax{} and $R^{2}$ from
\RelationLooRsqMin{} to \RelationLooRsqMax{}. The $\omega$-conotoxin point was displayed as
an external comparison and was not included in the Pfam fit. These data show a preliminary
association, not a prospective threshold or decision rule. The attention weights tracked
$f_{\mathrm{eff}}$ closely in all six families. Thus, the multiplicity-weighted energy reached
its target latent mass despite the downstream calibration gap. Hard curation achieved
100\% marker retention in every family, independent of the separation index. This result is
expected because its PCA basis contains only the designated subset.

We also tested whether increasing $\beta$ beyond $\beta^{*}$ transferred more marker signal
through the PCA representation. We swept $\beta/\beta^{*}$ from 0.5 to 3.0 at
$\rho = 10$, 50, and 200 on the Kunitz domain. Higher $\beta$ increased the P1~K/R
fraction and reduced diversity. At $\rho = 200$ and $\beta = 3\beta^{*}$, P1~K/R
reached 0.71, compared with 0.61 at $\beta^{*}$. Diversity decreased from 0.53 to 0.46.
Thus, $\rho$ and $\beta/\beta^{*}$ provide separate controls for marker retention and diversity.

\paragraph{The hard mask as the $\rho \to \infty$ endpoint.}
Hard attention masking sets the background attention mass to zero. It is the
$\rho \to \infty$ endpoint of the multiplicity axis on the full-family basis. It also gives
$\Delta_{\mathrm{attn}} = 0$ in Eq.~\ref{eq:gap-decomp}
(Table~\ref{tab:mask-recovery}).

We held the mask fixed and swept $\beta$ to measure the decode-limited recovery ceiling.
The P1~K/R fraction increased from 0.46 at $\beta = 2$ to 1.00 at $\beta = 512$
without a plateau (Fig.~\ref{fig:mask-betasweep}). Thus, the residual decode gap
$\Delta_{\mathrm{PCA}} + \Delta_{\mathrm{argmax}}$ decreased as retrieval sharpened.
At matched $\beta$, hard masking recovered more marker signal than finite-$\rho$ weighting.
The difference was 0.085 (SE 0.028) at $f_{\mathrm{eff}} = 0.5$ and 0.053
(SE 0.029) at $f_{\mathrm{eff}} = 0.7$. It decreased to 0.003 (SE 0.030) at
$f_{\mathrm{eff}} = 0.99$, where the mask and strongest finite-$\rho$ condition coincided.

Hard masking removes the background components and raises the intended designated mass from
$f_{\mathrm{eff}}$ to one. Matched-$\beta$ comparisons therefore have different intended
masses. Their difference should decrease as $f_{\mathrm{eff}}$ approaches one. Attention
realized the intended mass in each condition to within a fraction of a percent
($\Delta_{\mathrm{attn}} \approx 0$). A 0.5 difference in intended mass produced a
0.085 difference (SE 0.028) in observed marker fraction. The latent-to-sequence map
attenuated the response through noise, PCA reconstruction, and argmax decoding.
Both retrieval sharpness $\beta$ and conditioning strength $\rho$ affected recovery.
At high $\beta$, novelty decreased from 0.42 at the mask's $\beta^{*}$ to 0.16 at
$\beta = 512$ (Table~\ref{tab:mask-recovery}). Hard curation reached complete marker
retention at low $\beta$ with novelty of 0.32. Its subset-only basis retains designated
variation that the leading full-family components blur.

\paragraph{Application to a Cav2.2-associated $\omega$-conotoxin subset.}
We used the $\omega$-conotoxin O-superfamily to test curated conditioning beyond classical
protein domains (Table~\ref{tab:conotoxin-pharmacophore}, Fig.~\ref{fig:conotoxin-loop}).
These peptides come from cone snails and are rich in disulfide bonds. They block
voltage-gated calcium channels, including Cav2.2 in sensory neurons
~\cite{omegaConotoxinSAR2006}. Cav2.2 contributes to pain
signaling between peripheral nociceptors and the spinal cord. The family includes
$\omega$-conotoxin MVIIA (ziconotide), an approved non-opioid analgesic for intractable
pain~\cite{ziconotideReview2004}. Prior modeling of ATP-gated calcium dynamics in dorsal root
ganglion neurons also identified Cav2.2 as part of the pain signaling network
~\cite{songModelingAnalysis2009}.

Tyr13 is a reported Cav2.2 marker in MVIIA numbering. Its hydroxyl group contacts the channel
pore, and substitution of this residue abolishes binding. Basic residues distributed across
the peptide also contribute to electrostatic complementarity with the channel entrance.
We compiled 74 SwissProt $\omega$-conotoxin sequences and aligned them with
MAFFT~\cite{mafft2013}. Gap filtering retained 26 positions. The study's curated
Cav2.2-associated input list designated 23 accessions. The remaining 51 sequences
form a background set for conditioning and are not asserted to be nonbinders. An accession-level
audit records which designations have activity metadata in the repository and which are supported
only by the input list.

We ran two SA experiments with 50 chains, 5\,000 steps, and $\alpha = 0.01$. Each
condition generated 1\,550 sequences. One used the full family
($K = 74$, $\beta^{*} = 3.23$), and the other used only the designated accessions
($K = 23$, $\beta^{*} = 2.72$). Tyr occurred at the marker position in 33.8\% of the
full-family input and 46.9\% of full-family generation. It occurred in 82.6\% of the
designated input and 98.3\% of designated-subset generation
(Table~\ref{tab:conotoxin-pharmacophore}).

Accession designation and Tyr13 are distinct labels. They agree for 64 of 74 stored
sequences (86.5\%). The designated set contains 19 marker-positive sequences, and the
background contains six. We therefore report designation conditioning and Tyr13 recovery
separately. Designated-subset generation contained 20.1\% Lys or Arg, compared with
19.7\% in its input. Full-family generation contained 12.0\%, consistent with the lower
K/R content of the full-family input.

Per-position frequency heatmaps were similar for the designated input and its generated
sequences (Fig.~\ref{fig:conotoxin-loop}). Tyr13, the conserved Cys framework, and
variable-loop diversity were retained. The residual heatmaps showed small changes across
positions. Sequence-level analysis gave the same pattern
(Fig.~\ref{fig:conotoxin-sequence-analysis}). The five highest-pLDDT sequences retained
all six highly conserved positions, all six cysteines, and Tyr at the marker. They contained
5--11 substitutions at variable sites. The entropy correlation was $r = 0.745$ for
designated-subset generation and $r = 0.997$ for full-family generation. The lower
designated-subset value reflects reduced variation at co-varying loop positions.

We also compared 12 positions with published structure--activity relationship observations
(Table~\ref{tab:sar-agreement}). Designated-subset generation retained Tyr13 at 98.3\%.
It retained all six disulfide-forming cysteines at frequencies of at least 97.4\%. These
cysteines are invariant in the designated input. Designated-subset generation raised the
frequency of basic residues at loop~2 positions reported as important for channel
interaction (Arg10: 82.3\% K/R against 69.6\% in the designated input) and at
position~21 (66.2\% K/R against 47.8\%). It exceeded full-family generation at every
reported position outside the disulfide framework, where both conditions sat at or
near 1.00. The generated library therefore enriched several residues reported by
prior mutagenesis. These frequencies do not test mutation effects in the generated backgrounds
~\cite{kimTyr13Essential1995,satoBasicResidues1993,lewisConotoxinSAR2012}.
The top three designated-subset sequences had predicted disulfide-constrained loop
architectures consistent with an inhibitor cystine-knot-like fold
(Fig.~\ref{fig:fold-superposition}). Curated conditioning transferred sequence statistics
from the 23 designated accessions to the generated library. Sequence generation did not
require retraining or the Cav2.2 target structure.

\paragraph{AlphaFold2 multimer complex prediction.}
We compared complex-level predictions for SA-generated and natural $\omega$-conotoxins using
ColabFold AlphaFold2-multimer with the Cav2.2 pore-domain target
(Table~\ref{tab:docking-validation}). SA designated-subset ($n=10$; pTM
$0.221\pm0.007$, iPTM $0.105\pm0.008$, and interface pLDDT $37.2\pm5.4$).
The SA full-family values were $0.219\pm0.008$, $0.104\pm0.007$, and
$38.2\pm4.0$ ($n=10$). Natural controls had values of $0.228\pm0.004$,
$0.108\pm0.007$, and $40.7\pm3.5$ ($n=5$). All groups received low-confidence
complex predictions. Exact permutation tests detected no group
difference on the four reported metrics (all $p>0.05$), but these tests do not establish
equivalence or preserved binding geometry. Experimental complex structures and binding assays
are required to evaluate Cav2.2 interaction.

\paragraph{Model-based structure and sequence evaluation.}
We predicted folds for 50 sequences from each source (stored, full-family SA,
K/R-positive SA, K/R-negative SA, and HMM baseline) using ESMFold and computed TM-scores
against the experimental Kunitz domain structure (1BPI, chain~A) to compare
model-predicted structure summaries (Table~\ref{tab:structure-validation}).
SA-generated and stored sequences had similar model scores. Mean pLDDT was
$90.4 \pm 1.2$ for SA full-family, $90.7 \pm 1.1$ for SA K/R-positive, and
$89.3 \pm 3.2$ for stored sequences. All SA sequences exceeded the pLDDT $>70$
threshold. Their mean TM-scores were $0.84 \pm 0.01$, $0.83 \pm 0.01$, and
$0.83 \pm 0.03$, respectively (Fig.~\ref{fig:fold-superposition}). HMM-emitted
sequences had lower and more variable model scores. Their mean pLDDT was
$60.4 \pm 13.1$, with 20\% above the threshold, and their mean TM-score was
$0.60 \pm 0.19$ (Table~\ref{tab:structure-validation}).

AlphaFold2-ptm predictions via ColabFold gave similar summaries
~\cite{jumperHighlyAccurate2021,mirdita2022colabfold}. Kunitz TM-scores were
$0.83$--$0.85$ with pLDDT of $93$--$95$. The shorter $\omega$-conotoxins had
TM-scores of $0.35$--$0.47$ with pLDDT of $68$--$79$
(Table~\ref{tab:af2-validation}).

We used ESM2-650M masked marginal scoring to assess sequence plausibility. The resulting
pseudo-perplexity values were $4.78 \pm 0.64$ for SA full-family, $4.72 \pm 0.60$ for
SA K/R-positive, and $4.97 \pm 0.98$ for stored sequences. HMM-emitted sequences
had a value of $16.1 \pm 4.2$ (Table~\ref{tab:structure-validation}). The lower
SA values are consistent with generation near the center of the family distribution.

We also compared SA with five HMMER3 profile-HMM emissions of 150 sequences and five
bootstrap resamples~\cite{hmmer3}. The seed-42 HMM emission was also used for the model-based
evaluations. Full-family SA
gave a P1~K/R fraction of $0.38 \pm 0.03$ and KL divergence of
$0.0076 \pm 0.0009$. HMM emission gave $0.17 \pm 0.04$ and
$0.027 \pm 0.002$, respectively. Bootstrap resampling preserved composition
(KL $\approx 0.001$) but produced no novel sequences. K/R-positive-subset SA
gave a P1~K/R fraction of $1.0 \pm 0.0$. The unconditional baselines do not
provide this marker-conditioning control (Table~\ref{tab:structure-validation}).
